\documentclass[11pt,a4paper]{article}
\usepackage{fancyhdr}
\usepackage{notes}
\usepackage{amsmath}
\usepackage{enumitem}

\pagestyle{fancy}
\fancyhf{}
\rhead{Geophyics 3A: Week 4}
\rhead{Density and Porosity Activity}
\cfoot{\thepage}

\title{Worksheet A -- Density and Porosity Activiy}
\author{Geophysics 3A: Potential Fields \& Geothermics}
\date{\today}


\begin{document}
\section*{Introduction}
In week 2, we calculated several physical properties, though in each case we assumed porosity was negligible.  However, for many rocks porosity

Consider a sandstone composed of quartz ($\rho_q=2650$ kg m$^{-3}$) and feldspar ($\rho_f=2600$ kg m$^{-3}$).
The volumetric fractions of solids are $f_q$ and $f_f$ with $f_q + f_f = 1$.

\subsection*{Tasks}
\begin{enumerate}
  \item Write an expression for bulk density:
  %\[
  %  \rho_{\text{bulk}} = (1-\phi)(f_q\rho_q + f_f\rho_f) + \phi\rho_{\text{fluid}}.
  %\]

  \item Compute $\rho_b$ for:  
        (i) $\phi = 0.30$, air-filled pores ($\rho_{\text{fluid}}\approx 1$),  
        (ii) $\phi = 0.30$, water-filled pores ($\rho_{\text{fluid}}=1000$),  
        (iii) $\phi = 0.30$, brine-filled pores ($\rho_{\text{fluid}}=1200$).

  \item Repeat for $\phi = 0.10$ (burial compaction).

  \item Explain physically:  
        \begin{itemize}
          \item Why does $\rho_b$ increase with reduced $\phi$?  
          \item Which fluid substitution changes $\rho_b$ the most?  
        \end{itemize}
\end{enumerate}

\end{document}